\begin{frontmatter}

\title{Cheap Signals, Costly Proof: Award-Layer Triage in Cartel Enforcement\texorpdfstring{\footnote{This version: May 2026. We thank seminar participants at INSPER for comments on earlier drafts; all remaining errors are our own.}}{}}

\author[insper]{Darcio Genicolo-Martins\corref{cor1}}
\ead{darciogm1@insper.edu.br}
\author[insper]{Paulo Furquim de Azevedo}
\affiliation[insper]{organization={INSPER},
  city={S\~ao Paulo}, country={Brazil}}
\cortext[cor1]{Corresponding author.}

\begin{abstract}
Cartel cases begin under thin information. Agencies often see
award records long before they recover the bid-level evidence
needed for proof. We ask whether that cheap layer can guide where
costly forensics should start. Using São Paulo's BEC procurement
platform ($\valYearStart$--$\valYearEnd$), we build a
\emph{frequent-loser} screen from award records alone: among
zero-win firms, participation intensity ranks loser-side
cartel-adjacency risk. Computed from $2009$--$2016$
participation, the screen prospectively discriminates
adjudicated CADE cobidders in $2017$--$2019$. On the same target,
it matches a seven-feature Imhof--Wallimann bid-distribution
pipeline and adds non-redundant signal in combination. Used as a
gatekeeper, the award-layer screen cuts the bid-microdata pool
for the forensic stage by $\valArchFootprintReduction$ while
still flagging $\valArchTPThouSeqK$ of $\valArchNPos$
adjudicated cobidders. The statistic creates suspicion, not
proof: it ranks cartel-adjacent firms and environments for
forensic attention without adjudicating membership or estimating
damages.
\end{abstract}

\end{frontmatter}

\noindent\textbf{Keywords:} screening under incomplete observability, cover bidding, cartel adjacency, award-layer enforcement, separating equilibrium, participation-based statistics

\noindent\textbf{JEL No.:} D44, D73, H57, K21, L41

\bigskip
